% META: {"id":"1SPE-GEOREP-CO-037","chapitre":"1SPE-GEOMETRIE-REPEREE","type_objet":"corrige","exercice_ref":"1SPE-GEOREP-EX-037","capacites_codes":["C4"],"sources_inspiration":[],"status":"approved","origin":"generated","status_history":["generated","audited","accepted","approved"],"release_acceptance":"RELEASE_OWNER_BATCH_ACCEPTANCE","acceptance_closure_digest":"sha256:9b3ccf9a81c5520fbb7e03b7b2d3e7bf2057a02834b6d908bf3be5f49f3b553f","acceptance_receipt_digest":"sha256:4fad86f0282e0fa4e0419cca1ad6379ae7af5a280c04a4a90880f90a1c044242"}
% BEGIN-VERIFY
% from sympy import *
% # Tangence: cercle centre (2,-1) rayon sqrt(5), droite D: x - 2y + c = 0
% # d = |2+2+c|/sqrt(1+4) = |4+c|/sqrt(5) = sqrt(5)
% # |4+c| = 5 => c = 1 ou c = -9
% assert abs(4+1) == 5
% assert abs(4-9) == 5
% END-VERIFY
\begin{corrige}{1SPE-GEOREP-EX-037}

\commentaireMarge{Tangence $\iff$ $d(\Omega, \mathcal{D}) = r$.}

$d = \dfrac{|2 - 2(-1) + c|}{\sqrt{1 + 4}} = \dfrac{|4 + c|}{\sqrt{5}}$.

$d = \sqrt{5} \iff |4 + c| = 5$.

$4 + c = 5$ ou $4 + c = -5$, soit $\boxed{c = 1}$ ou $\boxed{c = -9}$.

\end{corrige}
